\documentclass[10pt]{article}

%% Esempio 2 - Verifica di fisica (intestazione compatta).
%% Mostra: materia fisica (colore verde), misure DSA, unita' con siunitx,
%% titolo lungo su una riga propria, diagrammi TikZ, vero/falso con choices,
%% tabella finale dei punteggi.
\usepackage[subject=physics]{fancytests}
\usepackage{lmodern}
\renewcommand*\familydefault{\sfdefault}
\usepackage[T1]{fontenc}

\geometry{a4paper, top=15mm, bottom=15mm, left=17.5mm, right=17.5mm}
\linespread{1.25}

\testsetup{class=4G, style=compact, dsa=rid}

\begin{document}

\testheader[date=22/05/2026,
  title={Onde meccaniche, suono e fenomeni ondulatori: battimenti ed effetto Doppler}]

\exercise Un'onda sinusoidale si propaga lungo una corda con velocità
$v = \qty{340}{m/s}$ e frequenza $f = \qty{170}{Hz}$.

\begin{parts}
  \itemx[pts=1] Calcola la lunghezza d'onda $\lambda$.
  \itemx[pts=1] Se la frequenza raddoppia a parità di mezzo, la lunghezza d'onda
    \begin{choices}
      \correct si dimezza.
      \choice  raddoppia.
      \choice  resta invariata.
      \choice  quadruplica.
    \end{choices}
  \itemx[pts=1, noreduced] Calcola il periodo dell'onda.
\end{parts}

\exercise[diff=2] Un blocco di massa $m = \qty{2.0}{kg}$ scivola senza attrito
lungo un piano inclinato di $30^\circ$ ed è collegato a una molla.

\begin{minipage}[c]{.42\linewidth}
  \centering
  \begin{tikzpicture}[scale=.9]
    \fill[ground] (0,0) -- (4.5,0) -- (4.5,2.6) -- cycle;
    \draw[thick] (0,0) -- (4.5,0) -- (4.5,2.6) -- cycle;
    \node[body, minimum size=7mm, rotate=30] (b) at (2.3,1.65) {};
    \draw[spring] (b.east) -- (4.5,2.6);
    \draw[force] (b.center) -- ++(0,-1.0) node[right] {$m\vb{g}$};
    \draw[velocity] (b.west) -- ++(-1.1,-0.63) node[above left] {$\vb{v}$};
    \draw (1.0,0) arc (0:30:1.0) node[midway, right] {\small $30^\circ$};
  \end{tikzpicture}
\end{minipage}\hfill
\begin{minipage}[c]{.55\linewidth}
  \begin{parts}[ragged]
    \itemx[pts=1] Disegna il diagramma di corpo libero del blocco.
    \itemx[pts=1.5] Calcola la componente della forza peso parallela al piano.
    \itemx[pts=1.5, noreduced] Determina l'allungamento della molla all'equilibrio,
      sapendo che $k = \qty{200}{N/m}$.
  \end{parts}
\end{minipage}

\exercise Stabilisci se le seguenti affermazioni sono vere o false.

\begin{parts}
  \itemx[pts=0.5] Il suono si propaga anche nel vuoto.
    \begin{choices}[shuffle=false] \choice Vero \correct Falso \end{choices}
  \itemx[pts=0.5] La frequenza dei battimenti è la differenza delle due frequenze.
    \begin{choices}[shuffle=false] \correct Vero \choice Falso \end{choices}
\end{parts}
\exercisescore

\vfill
\finalscore[style=table]

\end{document}
